Consider two census tracts separated by a single road in a mid-sized North Carolina city.
Both tracts are zoned within the same school district; their children attend the same
elementary and middle schools.
Both tracts are served by the same volunteer fire company, which they co-fund through
a special district property tax.
Their water and sewer service comes from the same municipal utility, and they purchase
electricity from the same electric cooperative at the same rates, subject to the same
state public utility commission.
In what sense are these two tracts \emph{not} members of the same community?

Administrative zone co-membership answers this question with precision.
The six major zone types that structure civic life in the United States---school
districts, county subdivisions, electric utility service territories, fire districts,
water and sewer districts, and police precincts---are not arbitrary bureaucratic
categories.
They represent the actual administrative units through which residents exercise
collective decisions about taxation, infrastructure, public safety, and education.
Two tracts that are co-members of many zone types share not just a neighborhood
but a civic fate: they vote on the same measures, receive the same services, and
bear the same tax burdens.
This is the strongest possible administrative definition of a community of interest.

\subsection{Redistricting and Administrative Zone Splitting}

When a redistricting plan splits a school district---placing some tracts in one
congressional or legislative district and other tracts of the same school district
in a different one---it separates neighbors who share a fiscal and civic unit.
Parents in the same school district who are represented by different legislators
have a harder time advocating for school funding as a unified bloc.
Fire district splitting creates the anomaly of a fire station that responds to
calls in two different legislative districts with potentially different levels of
fire code enforcement or public safety appropriations.

School district splitting is a recognized harm in redistricting litigation.
Expert witnesses and commission staff have cited school district integrity as a
community of interest criterion in proceedings across California, Colorado, Arizona,
Virginia, and North Carolina.
The Brennan Center's best-practices guide for redistricting commissions lists
school district preservation as a concrete operationalization of the communities
of interest criterion~\citep{levitt2011guide}.

Yet despite this recognition, no prior redistricting algorithm has incorporated
school district co-membership (or any other administrative zone type) directly
into its edge weight computation.
The \textsc{bisect} system's three-layer compositor provides the infrastructure
for exactly this capability: Layer 2 (edge weights) can be overridden with any
function of the tract adjacency graph, including the zone co-membership score
defined in this paper.

\subsection{Contribution}

This paper makes four contributions:
\begin{enumerate}
  \item \textbf{Zone co-membership formula}: We define the zone co-membership
    score and edge weight modifier precisely, with full specification of the
    graceful degradation behavior when zone data is partially unavailable
    (\S\ref{sec:formula}).
  \item \textbf{Data pipeline}: We identify and document the six data sources
    for zone co-membership in order of coverage priority, with URLs, formats,
    update frequencies, and availability statistics
    (\S\ref{sec:data-sources}).
  \item \textbf{Legal analysis}: We ground the zone co-membership criterion
    in the property tax fiscal bond doctrine and in the court cases that
    recognize administrative subdivision integrity as a legitimate redistricting
    consideration (\S\ref{sec:legal}).
  \item \textbf{Experimental design}: We specify the experiments that will
    validate the algorithm's empirical effect once implementation is complete,
    including NC-14 school district preservation, WI fire district clustering,
    and New England town integrity (\S\ref{sec:experimental-design}).
\end{enumerate}

\subsection{Relationship to M.0}

This paper is an instantiation of the community character framework defined in
M.0~\citep{m0framework}.
M.0 established the general form $w(u,v) = w_{\mathrm{base}} \times \mathrm{sim}(\mathbf{c}_u, \mathbf{c}_v)$.
M.6 replaces cosine similarity with the zone co-membership ratio formula, which
satisfies the same mathematical properties (symmetry, $[0,1]$ range, graceful
degradation) but is more directly interpretable for legal purposes: it counts
shared administrative relationships rather than computing a geometric angle between
character vectors.

\subsection{Paper Organization}

Section~\ref{sec:zones-as-community} motivates administrative zones as community
signals.
Section~\ref{sec:formula} defines the zone co-membership score and edge weight modifier.
Section~\ref{sec:data-sources} documents all data sources.
Section~\ref{sec:legal} provides the legal grounding.
Section~\ref{sec:experimental-design} describes the deferred empirical validation.
Section~\ref{sec:conclusion} concludes with an expert witness template.
